% Part: first-order-logic
% Chapter: semantics
% Section: intro-semantics

\documentclass[../../../include/open-logic-section]{subfiles}

\begin{document}

\olfileid{fol}{syn}{its}

\olsection{ভূমিকা}

রাশির অর্থ দেওয়া অর্থতত্ত্বের বিষয়। অর্থতত্ত্বের কেন্দ্রীয় ধারণা হলো
!!a{structure}-এ পরিতৃপ্তি। !!^a{structure} ভাষার নির্মাণ-উপাদানগুলিকে
অর্থ দেয়: !!a{domain} হলো বস্তুর একটি অশূন্য সেট। পরিমাণসূচকগুলিকে
এই সংজ্ঞাক্ষেত্রের উপর বিস্তৃত বলে ব্যাখ্যা করা হয়, !!{constant}s-কে
সংজ্ঞাক্ষেত্রের সদস্য আরোপ করা হয়, !!{function}s-কে !!{domain} থেকে
নিজের মধ্যেই ফলন এবং !!{predicate}s-কে !!{domain}-এর উপর সম্বন্ধ আরোপ
করা হয়। !!{domain} এবং মৌলিক শব্দভাণ্ডারে করা আরোপগুলি একত্রে
!!a{structure} গঠন করে। !!^{variable}s, !!{formula}s-এ থাকতে পারে; তাদের
অর্থতত্ত্ব দিতে !!{domain}-এর !!{element}s-ও তাদের আরোপ করতে হয়—এটিই
চলরাশি-আরোপ। শেষে পরিতৃপ্তি-সম্বন্ধ এই সবকিছুকে একত্র করে। কোনো
!!^a{formula}, !!a{structure}~$\Struct{M}$-এ !!a{variable} আরোপ~$s$-এর
সাপেক্ষে পরিতৃপ্ত হতে পারে; একে $\Sat{M}{!A}[s]$ লিখি। এই সম্বন্ধটিকেও
$!A$-এর গঠনের উপর আরোহ করে সংজ্ঞায়িত করা হয়। যেমন, $(!A \land !B)$-এর
পরিতৃপ্তিকে $!A$ ও~$!B$-এর পরিতৃপ্তি (বা অপরিতৃপ্তি) দিয়ে সংজ্ঞায়িত
করতে যৌক্তিক সংযোজকগুলির সত্যক-সারণি ব্যবহার করি। তারপর দেখা যায়,
!!{formula}~$!A$ যদি !!a{sentence} হয়, অর্থাৎ তার কোনো মুক্ত চলরাশি
না থাকে, তাহলে !!{variable} আরোপটি অপ্রাসঙ্গিক। ফলে কোনো
!!{sentence}s, !!{structure}s-এ কেবল পরিতৃপ্ত (বা অপরিতৃপ্ত)—এ কথাই
বলতে পারি।

!!{sentence}s-এর জন্য পরিতৃপ্তি-সম্বন্ধ $\Sat{M}{!A}$-এর ভিত্তিতে
তারপর বৈধতা, ফলশ্রুতি ও পরিতৃপ্তিযোগ্যতার মৌলিক অর্থতাত্ত্বিক ধারণাগুলি
সংজ্ঞায়িত করতে পারি। কোনো !!^a{sentence} বৈধ, $\Entails !A$, যদি
প্রত্যেক !!{structure} তাকে পরিতৃপ্ত করে। !!{sentence}s-এর কোনো সেট থেকে
তার ফলশ্রুতি হয়, $\Gamma \Entails !A$, যদি $\Gamma$-র
সব !!{sentence}s-কে পরিতৃপ্ত করে এমন প্রত্যেক !!{structure}, $!A$-কেও
পরিতৃপ্ত করে। আর !!{sentence}s-এর কোনো সেট পরিতৃপ্তিযোগ্য, যদি কোনো
!!{structure} সেটটির সব !!{sentence}s-কে একই সঙ্গে পরিতৃপ্ত করে।
!!{formula}s আরোহীভাবে সংজ্ঞায়িত এবং পরিতৃপ্তিও পালাক্রমে
!!{formula}s-এর গঠনের উপর আরোহ করে সংজ্ঞায়িত বলে আমাদের অর্থতত্ত্বের
ধর্মগুলি প্রমাণ করতে এবং সংজ্ঞায়িত অর্থতাত্ত্বিক ধারণাগুলির মধ্যে
সম্বন্ধ স্থাপন করতে আমরা আরোহ ব্যবহার করতে পারি।

\end{document}
